\documentclass{article}
\usepackage{amsmath, amssymb, amsthm}
\usepackage[margin=2.5cm]{geometry}

\theoremstyle{definition}
\newtheorem*{remark}{Remark}

\begin{document}

% ---------------------------------------------------------------------------
\section*{Two stochastic epidemic models}
% ---------------------------------------------------------------------------

We consider two discrete-time, branching-process models for the spread of
an infectious disease.  Time is measured in days.  Let $I_t \geq 0$ denote
the number of new infections (incidence) on day $t$, and let $R_t > 0$ be
the effective reproduction number on day $t$.

The \emph{serial interval weights} $w_s > 0$ ($s = 1, 2, \ldots$) satisfy
$\sum_{s=1}^\infty w_s = 1$ and encode the probability that a secondary
infection appears $s$ days after the primary case.  Write
\[
  F_r = \sum_{s=1}^r w_s
\]
for the cumulative serial interval distribution at lag $r$, so
$1 - F_r = \sum_{s=r+1}^\infty w_s$ is the probability that transmission
occurs after day $r$.

The \emph{dispersion parameter} $k > 0$ controls individual-level variation
in infectiousness.  Small $k$ corresponds to high heterogeneity
(superspreading); as $k \to \infty$ the negative-binomial offspring
distribution converges to a Poisson.

\begin{remark}
  Throughout, $\mathrm{NB}(\text{mean} = \mu, \text{disp} = k)$ denotes the
  negative-binomial distribution with probability mass function
  \[
    \mathbb{P}(X = n)
    = \binom{n + k - 1}{n}
    \left(\frac{k}{k + \mu}\right)^{\!k}
    \left(\frac{\mu}{k + \mu}\right)^{\!n},
    \quad n = 0, 1, 2, \ldots,
  \]
  mean $\mu$, variance $\mu + \mu^2/k$, and probability generating function
  (pgf) $G(s) = \bigl(k\,/\,(k + \mu(1-s))\bigr)^k$.
  Similarly, $\mathrm{Gamma}(\alpha, \beta)$ denotes the Gamma distribution
  with shape $\alpha$ and rate $\beta$, so density
  $f(x) = \beta^\alpha x^{\alpha-1} e^{-\beta x}/\Gamma(\alpha)$,
  mean $\alpha/\beta$, and variance $\alpha/\beta^2$.
\end{remark}

% ---------------------------------------------------------------------------
\section*{SSE model (Superspreading Events)}
% ---------------------------------------------------------------------------

In the \emph{superspreading-events} (SSE) model, the force of infection on
day $t$ is the weighted sum of past incidence,
\begin{equation}
  \lambda_t = \sum_{s=1}^{\infty} w_s\, I_{t-s},
\end{equation}
and new infections are negative-binomially distributed:
\begin{equation}
  I_t \mid \lambda_t \;\sim\;
  \mathrm{NB}\!\left(\text{mean} = R_t \lambda_t,\;\text{disp} =
  k\lambda_t\right).
\end{equation}
The conditional mean of $I_t$ is $R_t \lambda_t$ and its variance is
$R_t \lambda_t + R_t^2 \lambda_t / k$, so the NB inflates the Poisson
variance by the factor $1 + R_t \lambda_t / k$, capturing superspreading at
the population level.

% ---------------------------------------------------------------------------
\section*{SSI model (Superspreading Individuals)}
% ---------------------------------------------------------------------------

In the \emph{superspreading individuals} (SSI) model, heterogeneity
enters at the individual level.  Each infected individual has a
latent infectivity that is Gamma distributed with mean~1 and variance~$1/k$.

Each cohort of $I_t$ infected individuals then has an aggregate infectivity
\begin{equation}
  Y_t \mid I_t \;\sim\;
  \mathrm{Gamma}\!\left(\text{shape} = k I_t,\;\text{rate} = k\right),
\end{equation}
so $\mathbb{E}[Y_t \mid I_t] = I_t$ and
$\mathrm{Var}(Y_t \mid I_t) = I_t/k$.  New infections arise as a Poisson
process driven by past infectivity,
\begin{equation}
  I_t \mid (Y_s)_{s < t} \;\sim\;
  \mathrm{Poisson}\!\left(R_t \sum_{s=1}^{\infty} w_s\, Y_{t-s}\right).
\end{equation}

\begin{remark}
  Marginalising $Y_t$ from the SSI model recovers a negative-binomial
  offspring distribution (by the Poisson--Gamma mixture identity; see
  below), so the two models share the same mean--variance structure at the
  population level while differing in the \emph{mechanism} generating
  overdispersion.
\end{remark}

% ---------------------------------------------------------------------------
\section*{Extinction probability}
% ---------------------------------------------------------------------------

Henceforth assume the reproduction number is constant, $R_t = R_0$.  We
study the probability that the epidemic goes extinct, starting from a single
infectious case on day~0.

By branching-process theory, the extinction probability $q$ is the smallest
non-negative fixed point of the pgf of the offspring distribution.  For both
models the offspring distribution is
$\mathrm{NB}(\text{mean} = R_0, \text{disp} = k)$, whose pgf is
$G(s) = (k/(k + R_0(1-s)))^k$.  Setting $q = G(q)$ gives
\begin{equation}
  q = \left(1 + \frac{R_0}{k}(1 - q)\right)^{-k}.
\end{equation}
If $R_0 \leq 1$ the unique solution is $q = 1$ (extinction is certain); if
$R_0 > 1$ there is a unique solution $q \in (0,1)$, so the probability of a
major outbreak is $1 - q > 0$.

% ---------------------------------------------------------------------------
\subsection*{Extinction probability following 1 case on day~0,
then \(r\) days without cases}
% ---------------------------------------------------------------------------

Suppose we observe a single case on day~0 and no new cases on days
$1, \ldots, r$.  We compute the updated extinction probability $q_r$ for
the residual epidemic.

% ---- SSE ------------------------------------------------------------------
\subsubsection*{SSE model}

Each individual infected on day~0 generates cases at lag $s$ according to
an independent $\mathrm{NB}(\text{mean} = R_0 w_s, \text{disp} = k w_s)$
distribution.  Observing zero cases on days $1, \ldots, r$ implies that all
offspring at lags $s \leq r$ were zero, so only lags $s > r$ remain in play.

A crucial closure property of the NB family is that sums of independent NB
random variables with the \emph{same success probability}
$p = k_i/(k_i + \mu_i)$ are again NB.  Here every lag $s$ yields
$p = kw_s/(kw_s + R_0 w_s) = k/(k + R_0)$, which is independent of $s$.
Summing over all remaining lags, the index case's remaining offspring
distribution is
\begin{equation}
  \sum_{s=r+1}^{\infty} \mathrm{NB}(\text{mean} = R_0 w_s,\;
  \text{disp} = k w_s)
  = \mathrm{NB}\!\left(\text{mean} = R_0(1 - F_r),\;
  \text{disp} = k(1 - F_r)\right),
\end{equation}
with pgf $G_r(s) = \bigl(1 + (R_0/k)(1-s)\bigr)^{-k(1-F_r)} = G(s)^{1-F_r}$.
Each such remaining offspring, once born, initiates an independent epidemic
governed by the original NB($R_0$, $k$) offspring distribution, and so goes
extinct with the standard probability $q$ (the smallest fixed point of $G$).
The conditioning event constrains only the index case's offspring at lags
$\le r$ — it does not propagate to later generations — so
\begin{equation}
  q_r = G_r(q)
  = \left(1 + \frac{R_0}{k}(1 - q)\right)^{-k(1 - F_r)}
  = q^{\,1 - F_r}.
\end{equation}
As $r \to \infty$, $F_r \to 1$ and $q_r \to 1$: each additional day without
cases raises the extinction probability towards certainty.

% ---- SSI ------------------------------------------------------------------
\subsubsection*{SSI model}

It is more natural here to work with the latent infectivity of the index
case.  Let $y$ denote the realised infectivity of the single case on day~0,
with prior
\begin{equation}
  y \sim \mathrm{Gamma}(k,\, k), \qquad \mathbb{E}[y] = 1.
\end{equation}
Conditional on $y$, infections at each lag $s$ are independent with
\begin{equation}
  I_s \mid y \;\sim\; \mathrm{Poisson}(R_0\, w_s\, y).
\end{equation}
Since $y$ is unobserved, we update its distribution given the zero
observations before computing the offspring distribution.

\paragraph{Step 1: posterior update given no cases on days \(1,\dots,r\).}

By Bayes' theorem, with the $I_s$ conditionally independent Poisson given $y$,
\begin{align}
  f(y \mid I_1 = \cdots = I_r = 0)
  &\propto f(y)\;\mathbb{P}(I_1 = \cdots = I_r = 0 \mid y) \notag\\
  &= \frac{k^k}{\Gamma(k)}\, y^{k-1} e^{-ky}
  \prod_{s=1}^r e^{-R_0 w_s y} \notag\\
  &= \frac{k^k}{\Gamma(k)}\, y^{k-1}
  \exp\!\left(-y\Bigl[k + R_0 \sum_{s=1}^r w_s\Bigr]\right) \notag\\
  &= \frac{k^k}{\Gamma(k)}\, y^{k-1}\, e^{-(k + R_0 F_r)\,y}.
\end{align}
This is proportional to a $\mathrm{Gamma}(k,\; k + R_0 F_r)$ density, so
\begin{equation}
  y \mid (I_1 = \cdots = I_r = 0)
  \;\sim\; \mathrm{Gamma}(k,\; k + R_0 F_r).
\end{equation}
The posterior rate $k + R_0 F_r > k$ reflects the fact that observing no
early transmission is evidence that $y$ is small: the index case is likely
a low-infectivity individual.

\paragraph{Step 2: marginal offspring distribution for lags \(s > r\).}

Given $y$, the total remaining offspring is
\begin{equation}
  I_{\mathrm{rem}} \mid y
  \;=\; \sum_{s=r+1}^{\infty} I_s
  \;\sim\; \mathrm{Poisson}\!\left(R_0(1 - F_r)\,y\right).
\end{equation}
To marginalise over the posterior, set $\lambda = R_0(1 - F_r)\,y$.  Since
$y \sim \mathrm{Gamma}(k,\; k + R_0 F_r)$, the scaling property
($cX \sim \mathrm{Gamma}(\alpha, \beta/c)$ for $X \sim
\mathrm{Gamma}(\alpha, \beta)$)
gives
\begin{equation}
  \lambda \;\sim\;
  \mathrm{Gamma}\!\left(k,\;\frac{k + R_0 F_r}{R_0(1 - F_r)}\right).
\end{equation}
The \emph{Poisson--Gamma mixture identity} states that if
$X \mid \lambda \sim \mathrm{Poisson}(\lambda)$ and
$\lambda \sim \mathrm{Gamma}(\alpha, \beta)$, then marginally
$X \sim \mathrm{NB}(\text{size} = \alpha, \text{mean} = \alpha/\beta)$.
Applying this,
\begin{align}
  I_{\mathrm{rem}}
  &\;\sim\;
  \mathrm{NB}\!\left(\text{size} = k,\;
  \text{mean} = k \cdot \frac{R_0(1 - F_r)}{k + R_0 F_r}\right).
\end{align}
Compared to the unconditional mean $R_0(1-F_r)$, the effective mean is
reduced by the factor
\begin{equation}
  \frac{k}{k + R_0 F_r} \;<\; 1,
\end{equation}
which arises entirely from the Bayesian update: the zero streak tells us the
index case probably has low infectivity.

\paragraph{Step 3: extinction probability.}

Let $G_r$ denote the pgf of $I_{\mathrm{rem}}$. Each remaining offspring,
once born, initiates an independent epidemic with a fresh latent infectivity
drawn from the prior $\mathrm{Gamma}(k, k)$, so it is governed by the
original NB($R_0$, $k$) offspring distribution and goes extinct with the
standard probability $q$. The conditioning constrains only the index case's
own offspring counts, not those of later generations, hence
\begin{equation}
  q_r = G_r(q) =
  \left(1 + \frac{R_0(1 - F_r)}{k + R_0 F_r}(1 - q)\right)^{-k}.
\end{equation}

\end{document}
